\section{Discussion}

\subsection{Practical Implications}

For practitioners building autonomous agentic systems, the framework suggests several design heuristics, which we offer tentatively pending broader empirical validation. First, it may be advisable to begin with the simplest architecture that functions and to add roles only when bottleneck signals are clear and persistent; premature fission incurs coordination overhead without commensurate capability gain. Second, tracking the \emph{location} of bottlenecks over time may provide a more informative diagnostic of system health than monitoring any single performance metric. Third, preference comparison rather than scalar metrics may offer a more tractable and honest approach to architecture evaluation, given the multi-dimensional nature of system health. Fourth, designing for substrate independence (communicating via protocols rather than shared implementation details) enables component substitution, whether replacing one LLM with another or replacing an agent with a human, without disrupting organizational architecture. Fifth, the HITL boundary should be learnable rather than static, as human cognitive bandwidth is typically the scarcest resource in agentic deployments.

\subsection{Limitations}

This work has several significant limitations. The Wallfacer experiment is a single case study conducted with a single LLM substrate (Claude Code) on a single class of tasks (software engineering). The 51,000-line threshold at which self-diagnosis emerged may be model-specific and task-specific; we have no evidence that it represents a general phase transition. The fission principle, while consistent with our observations, has not been tested across different LLM architectures, task domains, or team configurations. The theoretical extensions in Sections 5 and 6 are speculative and should be understood as a research agenda rather than a set of validated claims. The formal connections to Lawvere's theorem, G\"odel's incompleteness, and dissipative structures are analogical rather than deductive, and it remains to be seen whether they can be made rigorous.

\subsection{Open Questions}

Several questions merit investigation. First, can a normalized complexity metric predict architectural fission thresholds across different substrates and task domains? Second, do autonomous agentic systems converge to stable configurations, oscillate, or drift unboundedly under sustained operation? The dissipative structure analogy suggests oscillation as the generic attractor, but this has not been empirically established. Third, when does role fission improve net system performance (capability gain exceeding coordination cost), and can this be connected to optimal stopping theory \cite{shiryaev2007}? Fourth, how should independent agentic ecosystems establish inter-system trust and negotiate shared protocols? MCP \cite{anthropic2024mcp} provides a communication primitive, but the trust and preference negotiation layers remain to be designed.
